% pst-numeric-doc.tex: manual of pst-numeric
% Build with tools/make-doc.sh (latex, dvips, ps2pdf).
% Copyright 2026 Buddy Ledger
% LaTeX Project Public License 1.3c or later, see LICENSE
\documentclass[11pt,a4paper]{article}
\usepackage[T1]{fontenc}
\usepackage{lmodern}
\usepackage[margin=2.5cm]{geometry}
\usepackage{amsmath,array,booktabs}
\usepackage{pst-numeric}
\usepackage{pst-plot}
\usepackage[hidelinks]{hyperref}

\newcommand\pkg{\texttt{pst-numeric}}
\newcommand\cs[1]{\texttt{\textbackslash#1}}
\newcommand\st[1]{\texttt{#1}}
% shows PostScript code and, under it, the result it leaves on the stack
\newcommand\demo[1]{\par\smallskip\noindent
  \texttt{\detokenize{#1}}\par\nopagebreak\noindent
  \hspace*{1.5em}$\to$\hspace{0.5em}%
  \begin{pspicture}(0,-0.05)(12,0.3)\psNumericShow{#1}\end{pspicture}\par\smallskip}
\newenvironment{reference}
  {\par\smallskip\noindent\begin{tabular}{@{}>{\raggedright\ttfamily}p{0.33\linewidth}>{\raggedright\arraybackslash}p{0.62\linewidth}@{}}\toprule}
  {\bottomrule\end{tabular}\par\medskip}

\title{\pkg: numbers with any number of digits, matrices, complex numbers
  and numerical methods for PSTricks\\[1ex]\large Version 1.2}
\author{Buddy Ledger}
\date{21 November 2026}

\setlength\emergencystretch{3em}
\begin{document}
\maketitle

\begin{abstract}
\pkg{} is a library of PostScript procedures for PSTricks. Its decimal
numbers (\emph{bignums}) carry as many significant digits as you ask for,
where PostScript's own reals have about seven. On top of that it has
matrices, complex numbers, roots of functions, integrals and polynomials,
each working on ordinary reals or on bignums. The calculations run in the
PostScript interpreter while the page is drawn, so they can feed plots
directly.
\end{abstract}

\section{Loading}

\begin{verbatim}
\usepackage{pst-numeric}            % every part
\usepackage[matrix,complex]{pst-numeric}
\end{verbatim}

The library has five parts, each in its own dvips header, so that a document
carries only what it uses. The options choose them; with no option, all five
are loaded.

\begin{reference}
\textrm{Option} & \textrm{Part}\\\midrule
bignum & decimal numbers with any number of digits (\st{pst-numeric-bignum.pro})\\
trig & CORDIC trigonometry on reals (\st{pst-numeric-trig.pro})\\
matrix & matrices of reals or bignums (\st{pst-numeric-matrix.pro})\\
complex & complex numbers; loads \st{matrix} too\\
calc & roots, integrals, polynomials, gcd; loads \st{matrix} and \st{complex} too\\
all & all five\\
\end{reference}

Matrices, complex numbers and \st{calc} work on reals without \st{bignum};
load it as well to use them with bignums. Instead of options, a document can
define \cs{pstNumericParts} before loading the package, as a list without
spaces: \verb|\def\pstNumericParts{bignum,matrix}|.

\pkg{} needs PSTricks and works with the \texttt{latex}--\texttt{dvips}
route (or \texttt{xelatex}), like every PSTricks package.

\section{Using the procedures}

Every procedure is defined in one PostScript dictionary,
\st{tx@NumericDict} (the macro \cs{pstNumericDict} expands to its name), so
that none of its names clash with PSTricks or with other packages. PostScript
code that uses the library must open that dictionary. \cs{pstNumeric} does
this: \verb|\pstNumeric{code}| expands to
\verb*| tx@NumericDict begin code end |, and can be used inside any
PostScript argument of PSTricks:

\begin{verbatim}
\psplot{0}{3}{\pstNumeric{[0.25 0 -1] x polyval}}
\end{verbatim}

\cs{psNumericShow}\verb|[|\textit{options}\verb|]{|\textit{code}\verb|}| runs
the code and prints the result it leaves on the stack, at the current point,
in \cs{psNumericFont} (\st{Courier}) at \cs{psNumericFontSize} (10\,pt). It
is a PSTricks object of zero width, so put it in a \st{pspicture} that leaves
room for the text. The options are PSTricks options, \st{linecolor} giving
the colour of the text. A bignum is printed with all its digits, a complex
number as $a + b\,$\st{i}, a matrix as \st{[a b; c d]}, a list as
\st{[a, b, c]}; a string is printed as it is. Every result in this manual
is printed this way, while the page is drawn.

A list of two numbers looks the same as a complex number, and a list of
complex numbers the same as a matrix. Write \st{pstnum\_list} after the code
to print such a result as a list:

\demo{[1 2 3] polyder pstnum_list}

\paragraph{Precision.} \st{n setbnsize} sets the number of significant
digits of bignums (20 at the start). Inside \cs{pstNumeric} it applies until
the end of the page, because dvips restores the PostScript memory after each
page; a later \st{setbnsize} changes it again. Anything else the code
defines with \st{def} also goes into \st{tx@NumericDict}, so avoid the
names of the library's procedures.

\section{Decimal numbers with any number of digits}

A bignum is made with \st{setbn} from an integer, a real or a string. Pass
strings for exact values: a real has only about seven digits, so
\st{0.1 setbn} is not exactly one tenth, but \st{(0.1) setbn} is.

\demo{(0.1) setbn (0.2) setbn addbn}
\demo{50 setbnsize pibn 20 setbnsize}
\demo{30 setbnsize 2 setbn sqrtbn 20 setbnsize}
\demo{1 setbn expbn}
\demo{(-50) setbn expbn}
\demo{2 setbn 90 setbn powbn}
\demo{25 factbn}
\demo{60 30 binombn}
\demo{(0.5) setbn asinbn 6 setbn mulbn}

\begin{reference}
\textrm{Procedure} & \textrm{Stack: inputs $\to$ outputs}\\\midrule
setbnsize & \st{n} $\to$ (digits of every later result)\\
setbn & integer, real, string or bignum $\to$ bignum\\
bncvs, bncvr & bignum $\to$ string, real\\
addbn subbn mulbn divbn & \st{a b} $\to$ bignum\\
negbn absbn truncbn roundbn & \st{a} $\to$ bignum\\
cmpbn & \st{a b} $\to$ $-1$, 0 or 1\\
eqbn nebn gtbn gebn ltbn lebn & \st{a b} $\to$ boolean (\st{a b gtbn} is $a>b$)\\
bniszero & \st{a} $\to$ boolean\\
sqrtbn expbn lnbn log10bn & \st{x} $\to$ bignum\\
powbn & \st{x y} $\to$ $x^y$ (exact for whole-number $y$)\\
sinbn cosbn tanbn & \st{x} $\to$ bignum ($x$ in radians)\\
atanbn asinbn acosbn & \st{x} $\to$ bignum\\
atan2bn & \st{y x} $\to$ angle of the point $(x,y)$, in $(-\pi,\pi]$\\
sinhbn coshbn tanhbn & \st{x} $\to$ bignum\\
factbn & \st{n} $\to$ $n!$\\
binombn & \st{n k} $\to$ $\binom{n}{k}$\\
pibn & $\to$ $\pi$\\
\end{reference}

\section{Matrices}

A matrix is made from its rows with \st{mfromrows}. If any element is a
bignum or a string, the calculation uses bignums. The determinant of
\st{[[0.1 0.2] [0.3 0.4]]} is $-0.02$, so the fourth and fifth examples
should give 0: with reals there is a small error, with strings there is
none.

\demo{[[1 2] [3 4]] mfromrows [[5 6] [7 8]] mfromrows mmul}
\demo{[[2 1] [1 3]] mfromrows mdet}
\demo{[[2 1] [1 3]] mfromrows minv}
\demo{[[0.1 0.2] [0.3 0.4]] mfromrows mdet 0.02 add}
\demo{[[(0.1) (0.2)] [(0.3) (0.4)]] mfromrows mdet (0.02) setbn addbn}
\demo{[[2 1] [1 3]] mfromrows [[3] [5]] mfromrows msolve}
\demo{[[1 1] [1 2] [1 3] [1 4]] mfromrows [[6] [5] [7] [10]] mfromrows mlstsq}
\demo{[[2 1] [1 2]] mfromrows meigsym pop pstnum_list}

\begin{reference}
\textrm{Procedure} & \textrm{Stack}\\\midrule
marray & \st{m n} $\to$ $m\times n$ zero matrix\\
mfromrows, mtorows & \st{[[row] [row] ...]} $\to$ \st{M}, and back\\
midentity & \st{n} $\to$ identity\\
mcopy mtranspose & \st{M} $\to$ matrix\\
mget, mput & \st{M i j} $\to$ value; \st{M i j value} $\to$ (changes \st{M})\\
msize & \st{M} $\to$ \st{m n}\\
madd msub mmul & \st{A B} $\to$ matrix\\
mscale & \st{M s} $\to$ $s\cdot M$\\
mdet minv & \st{A} $\to$ determinant, inverse\\
msolve & \st{A B} $\to$ \st{X} with $AX=B$\\
mlstsq & \st{A B} $\to$ least-squares \st{X} (Householder QR)\\
meigsym & \st{A} $\to$ \st{[values] V} for a symmetric \st{A} (Jacobi)\\
\end{reference}

\section{Complex numbers}

A complex number is the array \st{[re im]}; a plain number $x$ counts as
\st{[x 0]}. As with matrices, a bignum or string part makes the calculation
use bignums.

\demo{[1 2] [3 4] cmul}
\demo{[1 2] [3 4] cdiv}
\demo{[-4 0] csqrt}
\demo{[(0) 1] [0 1] cpow}
\demo{[(1) 2] 10 cpow}

\begin{reference}
\textrm{Procedure} & \textrm{Stack}\\\midrule
cpx & \st{re im} $\to$ \st{z}\\
cre cim cabs carg & \st{z} $\to$ real part, imaginary part, modulus, argument\\
cadd csub cmul cdiv & \st{a b} $\to$ \st{z}\\
cneg cconj csqrt cexp cln & \st{z} $\to$ \st{z}\\
cpow & \st{z w} $\to$ $z^w$\\
ccvs & \st{z} $\to$ string\\
\end{reference}

\section{Roots, integrals and polynomials}

A function is a procedure that takes $x$ and returns $f(x)$; with bignums it
receives a bignum and must return one. A polynomial is the list of its
coefficients, highest power first.

\demo{{ dup mul 2 sub } 0 2 froot}
\demo{{ dup mulbn 2 setbn subbn } (1) (2) froot}
\demo{{ dup mul } 0 1 fintegrate}
\demo{[1 0 -2] 3 polyval}
\demo{[1 2] [1 -2] polymul pstnum_list}
\demo{[1 -6 11 -6] polyroots pstnum_list}
\demo{[0 1 2 3 4] [1.1 1.9 5.2 9.8 17.1] 2 polyfit}
\demo{84 36 gcd}

\begin{reference}
\textrm{Procedure} & \textrm{Stack}\\\midrule
froot & \st{\{f\} a b} $\to$ root of $f$ between $a$ and $b$ (Illinois)\\
fintegrate & \st{\{f\} a b} $\to$ $\int_a^b f$ (adaptive Gauss--Legendre)\\
polyval & \st{p x} $\to$ $p(x)$, $x$ real, bignum or complex\\
polyder & \st{p} $\to$ $p'$\\
polyadd polymul & \st{p q} $\to$ polynomial\\
polyroots & \st{p} $\to$ every root, as complex numbers (Durand--Kerner)\\
polyfit & \st{[x ...] [y ...] d} $\to$ least-squares polynomial of degree $d$\\
gcd lcm & \st{a b} $\to$ integer\\
\end{reference}

\section{Plots}

The procedures work anywhere PSTricks takes PostScript code. Here the
parabola is fitted once, with \cs{pstVerb}, and its coefficients are kept in
\st{tx@NumericDict} as \st{fitcoef}; the second curve is the integral
$\int_0^x e^{-t^2}\,dt$, worked out by \st{fintegrate} for each point.

\begin{verbatim}
\pstVerb{\pstNumeric{/fitcoef [0 1 2 3 4] [1.1 1.9 5.2 9.8 17.1]
  2 polyfit def}}
\psplot{0}{4}{\pstNumeric{fitcoef x polyval}}
\psplot{0}{4}{\pstNumeric{{ dup mul neg 2.718281828 exch exp }
  0 x fintegrate 10 mul}}
\end{verbatim}

\begin{center}
\pstVerb{\pstNumeric{/fitcoef [0 1 2 3 4] [1.1 1.9 5.2 9.8 17.1] 2 polyfit def}}
\psset{xunit=2cm,yunit=0.3cm}
\begin{pspicture}(-0.3,-1.5)(4.4,19)
\psaxes[Dy=5,ticksize=0 4pt]{->}(0,0)(4.3,18.5)
\psdots[dotstyle=*](0,1.1)(1,1.9)(2,5.2)(3,9.8)(4,17.1)
\psplot[linecolor=red,linewidth=1pt]{0}{4}{\pstNumeric{fitcoef x polyval}}
\psplot[linecolor=blue,linewidth=1pt,plotpoints=60]{0}{4}{\pstNumeric{{ dup mul neg 2.718281828 exch exp } 0 x fintegrate 10 mul}}
\rput[r](3.3,15){\red fitted parabola}
\rput[l](0.5,11.5){\blue $10\int_0^x e^{-t^2}\,dt$}
\end{pspicture}
\end{center}

\section{Errors and limits}

Wrong input raises the usual PostScript errors: \st{rangecheck} (for
example the square root of a negative bignum, or \st{froot} without a sign
change), \st{undefinedresult} (division by zero, a singular matrix) and
\st{syntaxerror} (a string \st{setbn} cannot read). The error stops the
page, and Ghostscript or the PostScript viewer reports it.

Bignum arithmetic works one decimal digit at a time: instant at 20 digits,
slower at 100 (\st{pibn} about 5\,s the first time). Integrals with bignums
are slow; the first one at a given precision spends about 10\,s working out
the Gauss--Legendre points.

\section{More}

The library is also a standalone PostScript library for Ghostscript, with a
test suite of 581 tests and a showcase of every procedure, at
\url{https://github.com/buddyledger/postscript-math}. Its README describes
every procedure in more detail.

\pkg{} may be distributed and modified under the LaTeX Project Public
License 1.3c or later.

\end{document}
